% GENERATED FILE. Edit canonical lessons, then run npm run generate:books.
% canonical-source-hash: fnv1a64:a343fb6fc106218f
% canonical-lessons: ZH-C67-ma, ZH-W67-ma-horse, ZH-C67-gou, ZH-C67-mao, ZH-C67-niao, ZH-C67-yu

\chapter{Animals}
\label{ch:zh-animals}
\hlchaptermodality{67}

\begin{chapteropening}
\textbf{By the end of this chapter:} \emph{I can name horses, dogs, cats, birds and fish.}

Complete the last lesson of chapter 67: I can name horses, dogs, cats, birds and fish.
\end{chapteropening}

\section[\texorpdfstring{\zh{马} (mǎ)}{马 (mǎ)}]{\zh{马} (mǎ) — horse}

\label{lesson:ZH-C67-ma}



\begin{quote}
Before the new one: say the Mandarin for snow, then the Mandarin for wind.
\end{quote}

\subsection*{You'll want to know: \zh{马}}
\textbf{\zh{马}} — \emph{mǎ} — \textquotedblleft{}horse\textquotedblright{}.

\begin{quote}
\textbf{\zh{我喜欢白马}} \emph{wǒ xǐhuan bái mǎ} — I love white horses.
\end{quote}

The new character, in pieces: \zh{马} — a turn, a zigzag that hooks, and a bar along the bottom: the right-hand piece of \zh{吗} and \zh{妈}. The next lesson puts it in your hand.

\begin{cousinweb}
\textbf{\zh{马}} (\emph{mǎ}) is inherited from Middle Chinese \emph{maeX} (Guangyun, in Baxter's transcription); the dictionary explains its old character as the shape of a head, a mane, a tail and four legs.
\end{cousinweb}

\begin{grammarlens}[title={What you've built}]
One more word for everyday answers.
\end{grammarlens}

\subsection*{Your turn}
\begin{itemize}
\raggedright
  \item \emph{Say it:} \emph{mǎ}
  \item \emph{Say it:} \emph{mǎ}, once more
  \item \emph{Say it:} \emph{wǒ xǐhuan bái mǎ} — I love white horses.
  \item \emph{Recall:} say the Mandarin for snow, then the Mandarin for wind, then say \emph{mǎ} again
\end{itemize}

\begin{tcolorbox}[breakable,colback=teal!4,colframe=teal!35!black,title={Before you move on}]
What does \zh{马} mean? (\textquotedblleft{}Horse\textquotedblright{}.) Say it once more.
\end{tcolorbox}
\section[\texorpdfstring{\zh{马} (mǎ)}{马 (mǎ)}]{\zh{马} — a turn, a zigzag, a base}

\label{lesson:ZH-W67-ma-horse}



\begin{quote}
The word first, then the character inside it.

\begin{itemize}
\raggedright
  \item \emph{Recall:} say \emph{mǎ} — \textbf{R1}, one lesson back, before it decays
\end{itemize}

\begin{quote}
\zh{吗} \zh{妈}
\end{quote}

\emph{Read it:} each one aloud, with nothing said first — \emph{ma}, \emph{mā}
\end{quote}

\subsection*{Script — a turn, a zigzag, a base}
\hlblockfigure{\detokenize{figures/ZH-W67-ma-horse-filmstrip.pdf}}{How \zh{马} is written, stroke by stroke}

\begin{quote}
\zh{马}
\end{quote}

\textbf{One piece, three strokes.}

\textbf{A horizontal that turns down} at its right end.

\textbf{A vertical that turns right, then turns down again} and hooks to the left at the foot, all in one stroke.

\textbf{A horizontal} along the bottom, rising a little to the right.

\subsection*{Your turn}
\begin{itemize}
\raggedright
  \item \emph{Write it:} \zh{马}, all three, with the model in front of you
  \item \emph{Cover:} write it again from memory, and count the pieces, not the strokes
\end{itemize}

\begin{tcolorbox}[breakable,colback=teal!4,colframe=teal!35!black,title={Before you move on}]
How many strokes in \zh{马}? (\textbf{Three}.) Which stroke turns without lifting? (\textbf{The second}.) Which characters you know have it on the right? (\textbf{\zh{吗} and \zh{妈}}.)

Source: \href{https://raw.githubusercontent.com/chanind/hanzi-writer-data/68d10a4b21150cae5e1ebbd223eed289cf32d90c/data/\%E9\%A9\%AC.json}{Hanzi Writer Data \zh{马}, pinned stroke paths and medians}.
\end{tcolorbox}
\section[gǒu]{gǒu — dog}

\label{lesson:ZH-C67-gou}



\begin{quote}
Before the new one: say the Mandarin for wind, then the Mandarin for horse.
\end{quote}

\subsection*{You'll want to know: gǒu}
\emph{gǒu} — \textquotedblleft{}dog\textquotedblright{}.

\begin{quote}
\emph{nǐ xǐhuan gǒu ma} — Do you like dogs?
\end{quote}

Its characters come later; for now it is a word for the ear.

\begin{cousinweb}
\emph{gǒu} is inherited from Middle Chinese \emph{kuwX} (Guangyun, in Baxter's transcription), glossed 'dog'.
\end{cousinweb}

\begin{grammarlens}[title={What you've built}]
One more word for everyday answers.
\end{grammarlens}

\subsection*{Your turn}
\begin{itemize}
\raggedright
  \item \emph{Say it:} \emph{gǒu}
  \item \emph{Say it:} \emph{gǒu}, once more
  \item \emph{Say it:} \emph{nǐ xǐhuan gǒu ma} — Do you like dogs?
  \item \emph{Recall:} say the Mandarin for wind, then the Mandarin for horse, then say \emph{gǒu} again
\end{itemize}

\begin{tcolorbox}[breakable,colback=teal!4,colframe=teal!35!black,title={Before you move on}]
What does gǒu mean? (\textquotedblleft{}Dog\textquotedblright{}.) Say it once more.
\end{tcolorbox}
\section[māo]{māo — cat}

\label{lesson:ZH-C67-mao}



\begin{quote}
Before the new one: say the Mandarin for horse, then the Mandarin for dog.
\end{quote}

\subsection*{You'll want to know: māo}
\emph{māo} — \textquotedblleft{}cat\textquotedblright{}.

\begin{quote}
\emph{wǒ jiā yǒu māo} — We've got a cat at home.
\end{quote}

Its characters come later; for now it is a word for the ear.

\begin{cousinweb}
\emph{māo} is inherited from Middle Chinese \emph{mjew} (Guangyun, in Baxter's transcription), glossed 'the animal that catches mice', though māo is not the sound \emph{mjew} would regularly give.
\end{cousinweb}

\begin{grammarlens}[title={What you've built}]
One more word for everyday answers.
\end{grammarlens}

\subsection*{Your turn}
\begin{itemize}
\raggedright
  \item \emph{Say it:} \emph{māo}
  \item \emph{Say it:} \emph{māo}, once more
  \item \emph{Say it:} \emph{wǒ jiā yǒu māo} — We've got a cat at home.
  \item \emph{Recall:} say the Mandarin for horse, then the Mandarin for dog, then say \emph{māo} again
\end{itemize}

\begin{tcolorbox}[breakable,colback=teal!4,colframe=teal!35!black,title={Before you move on}]
What does māo mean? (\textquotedblleft{}Cat\textquotedblright{}.) Say it once more.
\end{tcolorbox}
\section[niǎo]{niǎo — bird}

\label{lesson:ZH-C67-niao}



\begin{quote}
Before the new one: say the Mandarin for dog, then the Mandarin for cat.
\end{quote}

\subsection*{You'll want to know: niǎo}
\emph{niǎo} — \textquotedblleft{}bird\textquotedblright{}.

\begin{quote}
\emph{wǒ xǐhuan kàn niǎo} — I love watching birds.
\end{quote}

Its characters come later; for now it is a word for the ear.

\begin{cousinweb}
\emph{niǎo} is inherited from Middle Chinese \emph{tewX} (Guangyun, in Baxter's transcription), glossed, through an older book, as the name for all long-tailed birds. Why the modern word begins with n is not settled.
\end{cousinweb}

\begin{grammarlens}[title={What you've built}]
One more word for everyday answers.
\end{grammarlens}

\subsection*{Your turn}
\begin{itemize}
\raggedright
  \item \emph{Say it:} \emph{niǎo}
  \item \emph{Say it:} \emph{niǎo}, once more
  \item \emph{Say it:} \emph{wǒ xǐhuan kàn niǎo} — I love watching birds.
  \item \emph{Recall:} say the Mandarin for dog, then the Mandarin for cat, then say \emph{niǎo} again
\end{itemize}

\begin{tcolorbox}[breakable,colback=teal!4,colframe=teal!35!black,title={Before you move on}]
What does niǎo mean? (\textquotedblleft{}Bird\textquotedblright{}.) Say it once more.
\end{tcolorbox}
\section[yú]{yú — fish}

\label{lesson:ZH-C67-yu}



\begin{quote}
Before the new one: say the Mandarin for cat, then the Mandarin for bird.
\end{quote}

\subsection*{You'll want to know: yú}
\emph{yú} — \textquotedblleft{}fish\textquotedblright{}.

\begin{quote}
\emph{wǒ xǐhuan chī yú} — I love eating fish.
\end{quote}

Its characters come later; for now it is a word for the ear.

\begin{cousinweb}
\emph{yú} is inherited from Middle Chinese \emph{ngjo} (Guangyun, in Baxter's transcription), glossed, through an older book, as a water creature.
\end{cousinweb}

\begin{grammarlens}[title={What you've built}]
One more word for everyday answers.
\end{grammarlens}

\subsection*{Your turn}
\begin{itemize}
\raggedright
  \item \emph{Say it:} \emph{yú}
  \item \emph{Say it:} \emph{yú}, once more
  \item \emph{Say it:} \emph{wǒ xǐhuan chī yú} — I love eating fish.
  \item \emph{Recall:} say the Mandarin for cat, then the Mandarin for bird, then say \emph{yú} again
\end{itemize}

\begin{tcolorbox}[breakable,colback=teal!4,colframe=teal!35!black,title={Before you move on}]
What does yú mean? (\textquotedblleft{}Fish\textquotedblright{}.) Say it once more.
\end{tcolorbox}
